% block: matrix-equations-heading page:75 type:section
\subsection{{Matrix Equations}}\label{sec:matrix-equations}
{Consider the following matrix equation:}
\begin{equation}\label{eq:sylvester-equation}
AX+XB=C,
\end{equation}
{where $A$ and $B$ are $m$-by-$m$ and $n$-by-$n$ matrices. The matrix equation is known as the Sylvester equation, which arises in the field of control theory. The Sylvester equation also arises in a derivation of Newton's method for computing the matrix square root, see Sect.\ \ref{subsec:newton-matrix-square-root}.}
{The Sylvester equation can be transformed into linear systems of the form $M\boldsymbol{x}=\boldsymbol{c}$, and in what follows, we will see the connection between the Sylvester equation in \eqref{eq:sylvester-equation} and linear systems $M\boldsymbol{x}=\boldsymbol{c}$.}
% block: vectorization-definition page:75 type:ordinary-paragraph
{Let $G$ be an $n$-by-$m$ matrix whose column vectors are written as $\boldsymbol{g}_1,\boldsymbol{g}_2,\ldots,\boldsymbol{g}_m$. Then, matrix $G$ is written as $G=[\boldsymbol{g}_1,\boldsymbol{g}_2,\ldots,\boldsymbol{g}_m]$. We now define vec operator. The vec operator is a map from $\mathbb{C}^{n\times m}$ to $\mathbb{C}^{nm}$, and $\operatorname{vec}(G)$ is written as}
\begin{equation}\label{eq:vectorization-definition}
\operatorname{vec}(G)=\begin{bmatrix}\boldsymbol{g}_1\\\vdots\\\boldsymbol{g}_m\end{bmatrix}.
\end{equation}
{Some properties of the vec operator are listed next.}
% block: vectorization-properties page:76 type:theorem
\begin{Theorem}\label{thm:vectorization-properties}
{Let $X\in\mathbb{C}^{m\times n}$, $A_i\in\mathbb{C}^{m\times m}$, $B_i\in\mathbb{C}^{n\times n}$, and $P_{mn}$ be a permutation matrix given in Theorem\ \ref{thm:kronecker-permutation}. Then,}
\begin{enumerate}
\item $\operatorname{vec}(A_1XB_1+\cdots+A_kXB_k)=(B_1^{\top}\otimes A_1+\cdots+B_k^{\top}\otimes A_k)\operatorname{vec}(X)$,
\item $\operatorname{vec}(X^{\top})=P_{mn}\operatorname{vec}(X)$.
\end{enumerate}
\end{Theorem}
% block: vectorization-transpose-proof page:76 type:proof
\begin{Proof}
{We first give a proof of (2). From the definition of $P_{mn}$ in Theorem\ \ref{thm:kronecker-permutation},}
\[P_{mn}\operatorname{vec}(X)=\begin{bmatrix}
(E_{11})^{\top}&\cdots&(E_{1n})^{\top}\\\vdots&\ddots&\vdots\\(E_{m1})^{\top}&\cdots&(E_{mn})^{\top}
\end{bmatrix}\begin{bmatrix}\boldsymbol{x}_1\\\vdots\\\boldsymbol{x}_n\end{bmatrix},\]
{and thus the $k$th block element of $P_{mn}\operatorname{vec}(X)$ is given by}
\[\begin{aligned}
(P_{mn}\operatorname{vec}(X))_k&=\sum_{j=1}^nE_{kj}^{\top}\boldsymbol{x}_j
=\sum_{j=1}^nE_{jk}\boldsymbol{x}_j=\sum_{j=1}^n\boldsymbol{e}_j\boldsymbol{e}_k^{\top}\boldsymbol{x}_j
=\sum_{j=1}^n\boldsymbol{e}_jx_{kj}=\begin{bmatrix}x_{k1}\\\vdots\\x_{kn}\end{bmatrix},
\end{aligned}\]
{which is the $k$th column of $X^{\top}$. Therefore,}
\[P_{mn}\operatorname{vec}(X)=\begin{bmatrix}(P_{mn}\operatorname{vec}(X))_1\\\vdots\\(P_{mn}\operatorname{vec}(X))_n\end{bmatrix}=\operatorname{vec}(X^{\top}).\]
{Next, we give a proof of (1). Since $\operatorname{vec}(X+Y)=\operatorname{vec}(X)+\operatorname{vec}(Y)$, it is sufficient to prove $\operatorname{vec}(AXB)=(B^{\top}\otimes A)\operatorname{vec}(X)$. Let $M=XB$ and $M=[\boldsymbol{m}_1,\ldots,\boldsymbol{m}_n]$, then $AM=[A\boldsymbol{m}_1,\ldots,A\boldsymbol{m}_n]$. Thus}
\begin{equation}\label{eq:vectorized-left-multiplication}
\operatorname{vec}(AXB)=\operatorname{vec}(AM)=\begin{bmatrix}A\boldsymbol{m}_1\\\vdots\\A\boldsymbol{m}_n\end{bmatrix}
=(I\otimes A)\begin{bmatrix}\boldsymbol{m}_1\\\vdots\\\boldsymbol{m}_n\end{bmatrix}=(I_n\otimes A)\operatorname{vec}(M),
\end{equation}
{and the $\operatorname{vec}(M)$ is calculated as}
\[\begin{aligned}
\operatorname{vec}(XB)&=\operatorname{vec}((B^{\top}X^{\top})^{\top})=P_{nm}\operatorname{vec}(B^{\top}X^{\top})=P_{nm}(I_m\otimes B^{\top})\operatorname{vec}(X^{\top})\\
&=P_{nm}(I_m\otimes B^{\top})P_{mn}\operatorname{vec}(X)
=P_{mn}^{\top}(I_m\otimes B^{\top})P_{mn}\operatorname{vec}(X)=(B^{\top}\otimes I_m)\operatorname{vec}(X).
\end{aligned}\]
{The last equation holds from \eqref{eq:kronecker-swapping-permutation}, and this result together with \eqref{eq:vectorized-left-multiplication} yields}
\[\operatorname{vec}(AXB)=(I_n\otimes A)\operatorname{vec}(XB)=(I_n\otimes A)(B^{\top}\otimes I_m)\operatorname{vec}(X)=(B^{\top}\otimes A)\operatorname{vec}(X).\]
{The last equation follows from Theorem\ \ref{thm:kronecker-product-properties}-(6), which concludes the proof.}
\end{Proof}
% block: sylvester-linear-system page:77 type:ordinary-paragraph
{We are now ready to describe a relation between the Sylvester equation and linear systems. Applying Theorem\ \ref{thm:vectorization-properties}-(1) to Sylvester Eq.\ \eqref{eq:sylvester-equation} yields}
\[\operatorname{vec}(AX+XB)=\operatorname{vec}(C)\Leftrightarrow
\underbrace{[(I_n\otimes A)+(B^{\top}\otimes I_m)]}_M
\underbrace{\operatorname{vec}(X)}_{\boldsymbol{x}}=\underbrace{\operatorname{vec}(C)}_{\boldsymbol{c}},\]
{where $I_m$ is the $m$-by-$m$ identity matrix. Here we consider a slightly modified matrix equation of the form}
\[AX+X^{\top}B=C.\]
{The matrix is referred to as the T-congruence Sylvester equation. Applying the vec operator and using Theorem\ \ref{thm:vectorization-properties}-(2) yields}
\[\begin{aligned}
AX+X^{\top}B=C&\Leftrightarrow\operatorname{vec}(AX+X^{\top}B)=\operatorname{vec}(C)\\
&\Leftrightarrow\operatorname{vec}(AX)+\operatorname{vec}(X^{\top}B)=\operatorname{vec}(C)\\
&\Leftrightarrow\operatorname{vec}(I_n\otimes A)\operatorname{vec}(X)+(B^{\top}\otimes I_m)\operatorname{vec}(X^{\top})=\operatorname{vec}(C)\\
&\Leftrightarrow(I_n\otimes A)\operatorname{vec}(X)+(B^{\top}\otimes I_m)P_{mn}\operatorname{vec}(X)=\operatorname{vec}(C)\\
&\Leftrightarrow\underbrace{[(I_n\otimes A)+(B^{\top}\otimes I_m)P_{mn}]}_{\widetilde M}\underbrace{\operatorname{vec}(X)}_{\boldsymbol{x}}=\underbrace{\operatorname{vec}(C)}_{\boldsymbol{c}}.
\end{aligned}\]
{Therefore, the T-congruence Sylvester equation is also equivalent to linear systems of the form $\widetilde M\boldsymbol{x}=\boldsymbol{c}$.}
